\chapter*{新增参考文献}
\addcontentsline{toc}{chapter}{新增参考文献}
\SourcePage{191}
\begin{description}[style=nextline,leftmargin=5em,labelwidth=4.5em,font=\normalfont,itemsep=.7\baselineskip]
\item[{[BS15]}] B. Bhatt, P. Scholze，\term{scheme}{概形}的 pro-\term{etale}{平展}拓扑（\emph{The pro-étale topology for schemes}），载 \emph{De la Géométrie algébrique aux formes automorphes (I)}（从代数几何到自守形式，第 I 卷；纪念 Gérard Laumon 六十岁生日的论文集），Astérisque 369 (2015)，第 99--201 页。
\item[{[dC15]}] M. de Cataldo，有限域上的\term{proper}{固有}环面态射（\emph{Proper toric maps over finite fields}），Int. Math. Res. Not. IMRN 2015，第 24 期，第 13106--13121 页。
\item[{[Eke90]}] T. Ekedahl，论进形式体系（\emph{On the adic formalism}），载 \emph{The Grothendieck Festschrift}，第 II 卷，Progr. Math. 87 (1990)，Birkhäuser，第 197--218 页。
\item[{[Gab04]}] O. Gabber，关于若干 \term{t-structure}{$t$-结构}的札记（\emph{Notes on some $t$-structures}），载 \emph{Geometric aspects of Dwork theory}，第 II 卷 (2004)，Walter de Gruyter，第 711--734 页。
\item[{[GL17]}] D. Gaitsgory, J. Lurie，函数域上的 Weil 猜想 I（\emph{Weil's Conjecture for Function Fields I}），预印本，2017。
\item[{[Ham95]}] H. Hamm，\term{affine}{仿射}簇与 Lefschetz 定理（\emph{Affine varieties and Lefschetz theorems}），载 \emph{Singularity theory}（Trieste, 1991），第 248--262 页，World Sci. Publ., River Edge, NJ, 1995。
\item[{[Ill94]}] L. Illusie，围绕局部单值定理（\emph{Autour du théorème de monodromie locale}），载 \emph{Périodes $p$-adiques}（Bures-sur-Yvette, 1988），Astérisque 223 (1994)，第 9--57 页。
\item[{[Ill03]}] L. Illusie，\term{perversity}{反常度}与变动（\emph{Perversité et variation}），Manuscripta Math. 112 (2003)，第 271--295 页。
\item[{[Kas75]}] M. Kashiwara，极大超定线性微分方程组 I（\emph{On the maximally overdetermined system of linear differential equations. I}），Publ. Res. Inst. Math. Sci. 10 (1974/75)，第 563--579 页。
\item[{[KS90]}] M. Kashiwara, P. Schapira，流形上的\term{sheaf}{层}（\emph{Sheaves on manifolds}），Grundlehren der Mathematischen Wissenschaften 292 (1990)，Springer-Verlag。
\item[{[KL85]}] N. Katz, G. Laumon，\term{fourier-transform}{Fourier 变换}与指数和的上界（\emph{Transformation de Fourier et majoration de sommes exponentielles}），Publ. Math. IHÉS 62 (1985)，第 361--418 页。
\SourcePage{192}
\item[{[Lau83]}] G. Laumon，论\term{filtration}{滤过} $\mathcal D$-模的\term{derived-category}{导出范畴}（\emph{Sur la catégorie dérivée des $\mathcal D$-modules filtrés}），载 \emph{Algebraic geometry}（Tokyo/Kyoto, 1982），Lecture Notes in Math. 1016 (1983)，第 151--237 页。
\item[{[Lau87]}] G. Laumon，\term{fourier-transform}{Fourier 变换}、函数方程常数与 Weil 猜想（\emph{Transformation de Fourier, constantes d'équations fonctionnelles et conjecture de Weil}），Publ. Math. IHÉS 65 (1987)，第 131--210 页。
\item[{[MO15]}] D. Madore, F. Orgogozo，模 $\ell$ \term{etale}{平展}\term{cohomology}{上同调}的可计算性（\emph{Calculabilité de la cohomologie étale modulo $\ell$}），Algebra Number Theory 9 (2015)，第 7 期，第 1647--1739 页。
\item[{[Moc15]}] T. Mochizuki，\term{mixed}{混合}旋量 $\mathcal D$-模（\emph{Mixed twistor $\mathcal D$-modules}），Lecture Notes in Mathematics 2125 (2015)，Springer。
\item[{[PS14]}] V. Pilloni, B. Stroh，\term{affine}{仿射} Lefschetz 定理（\emph{Théorème de Lefschetz affine}），载 \emph{Travaux de Gabber sur l'uniformisation locale et la cohomologie étale des schémas quasi-excellents}（Gabber 关于局部一致化与拟优\term{scheme}{概形}\term{etale}{平展}\term{cohomology}{上同调}的工作），L. Illusie, Y. Laszlo, F. Orgogozo 编，Astérisque 363--364 (2014)，第 293--300 页。
\item[{[Sab13]}] C. Sabbah，野 Hodge 理论与 Hitchin--Kobayashi 对应（据 T. Mochizuki 的工作）（\emph{Théorie de Hodge et correspondance de Hitchin-Kobayashi sauvages (d'après T. Mochizuki)}），Séminaire Bourbaki 1050，Astérisque 352 (2013)，第 205--241 页。
\item[{[Sai88]}] M. Saito，可极化的 Hodge 模（\emph{Modules de Hodge polarisables}），Publ. Res. Inst. Math. Sci. 24 (1988)，第 849--995 页。
\item[{[Sai90]}] M. Saito，\term{mixed}{混合} Hodge 模（\emph{Mixed Hodge modules}），Publ. Res. Inst. Math. Sci. 26 (1990)，第 221--333 页。
\item[{[SZ16]}] S. Sun, W. Zheng，\term{intersection-cohomology}{交叉上同调}与通常\term{cohomology}{上同调}中的奇偶性和对称性（\emph{Parity and symmetry in intersection and ordinary cohomology}），Algebra Number Theory 10 (2016)，第 2 期，第 235--307 页。
\item[{[Wae72]}] L. Waelbroeck，$b$-空间的商（\emph{Les quotients de $b$-espaces}），载 \emph{Actes du Colloque d'Analyse Fonctionnelle}（Univ. Bordeaux, Bordeaux, 1971），Mém. Soc. Math. France，第 31--32 期 (1972)，第 389--394 页。
\end{description}
